\begin{frame}{Overview}
  \small
  \begin{multicols}{2}
    \begin{itemize}
      \item \textbf{1 — the foundations}
      \begin{itemize}
        \item interoperability: what is actually involved
        \item SGAM: where openADR sits — \alert{in a hands-on exercise}
        \item the Ship of Theseus: identity, meaning \& MDM
        \item information modelling 101: L0--L3, RDF, SKOS, OWL, SHACL
        \item the requirements — the yardstick the rest of the day is built on
      \end{itemize}
      \item \textbf{2 — the requirements, addressed one at a time}
      \begin{itemize}
        \item Req.\ 1: kinds of information models — vocabulary vs.\ constraint
        \item Req.\ 2: MOF — metamodels, expressivity, transformation
        \item meaning: OWA/CWA, the open world
        \item Req.\ 3: ontology alignment, upper ontologies
      \end{itemize}
      \item \textbf{3 — in practice, and the lens turned back on openADR}
      \begin{itemize}
        \item profiles: use-case specific models
        \item the toolchain: Avro, JSON Schema, SQL — \emph{generated}
        \item Netbeheer Nederland: the practice
        \item the three requirements, applied to openADR 3.0
      \end{itemize}
    \end{itemize}
  \end{multicols}
\end{frame}
